\documentclass[11pt,a4paper]{article}

\usepackage[utf8]{inputenc}
\usepackage[T1]{fontenc}
\usepackage{lmodern}
\usepackage{amsmath,amssymb,amsthm}
\usepackage{braket}
\usepackage{bm}
\usepackage{graphicx}
\usepackage{geometry}
\usepackage{xcolor}
\usepackage{listings}
\usepackage{booktabs}
\usepackage{hyperref}
\usepackage{enumitem}
\usepackage{microtype}
\usepackage{algorithm}
\usepackage{algpseudocode}

\geometry{margin=2.5cm}

\hypersetup{
  colorlinks=true,
  linkcolor=blue!60!black,
  citecolor=blue!60!black,
  urlcolor=blue!60!black,
}

% Code listings style
\definecolor{codegray}{rgb}{0.95,0.95,0.95}
\definecolor{codekeyword}{rgb}{0.0,0.3,0.7}
\definecolor{codecomment}{rgb}{0.3,0.5,0.3}
\definecolor{codestring}{rgb}{0.6,0.2,0.2}

\lstdefinelanguage{Fortran08}{
  morekeywords={subroutine,end,function,module,use,implicit,none,real,integer,logical,
                type,class,intent,in,out,inout,allocatable,allocate,deallocate,if,then,
                else,elseif,endif,do,while,call,return,parameter,contains,public,private,
                optional,dimension,target,pointer,associate,block,exit,cycle,select,case,
                true,false,dp,only},
  keywordstyle=\color{codekeyword}\bfseries,
  commentstyle=\color{codecomment}\itshape,
  stringstyle=\color{codestring},
  morecomment=[l]{!},
  morestring=[b]",
  morestring=[b]',
  sensitive=false,
}

\lstdefinestyle{hsd}{
  basicstyle=\ttfamily\small,
  backgroundcolor=\color{codegray},
  commentstyle=\color{codecomment}\itshape,
  morecomment=[l]{\#},
  frame=single,
  breaklines=true,
  columns=flexible,
}

\lstdefinestyle{fortran}{
  language=Fortran08,
  basicstyle=\ttfamily\small,
  backgroundcolor=\color{codegray},
  frame=single,
  breaklines=true,
  columns=flexible,
  numbers=left,
  numberstyle=\tiny\color{gray},
}

% Short commands
\newcommand{\kkappa}{\bm{\kappa}}
\newcommand{\xx}{\mathbf{x}}
\newcommand{\gg}{\mathbf{g}}
\newcommand{\HH}{\mathbf{H}}
\newcommand{\FF}{\mathbf{F}}
\newcommand{\CC}{\mathbf{C}}
\newcommand{\UU}{\mathbf{U}}
\newcommand{\Pp}{\mathbf{P}}
\newcommand{\II}{\mathbf{I}}
\newcommand{\SS}{\mathbf{S}}
\newcommand{\qq}{\mathbf{q}}
\DeclareMathOperator{\tr}{tr}

\title{\textbf{Second-Order SCF (SOSCF) for Excited-State Orbital\\
Optimisation in DFTB+:\\ Implementation Notes}}
\author{}
\date{\today}

\begin{document}

\maketitle

\begin{abstract}
This document describes a second-order SCF (SOSCF) implementation for
non-aufbau excited-state optimisation added to DFTB+ 24.1. The algorithm
combines (i) a level-shifted \emph{pre-SOSCF} Delta-SCF loop (STEP) that
drives the orbital gradient below a handoff threshold, with (ii) a
quasi-Newton orbital-rotation outer loop using the L-SR1 Hessian update
and the Cayley approximant of the matrix exponential.
Initial Maximum Overlap Method (IMOM) tracks orbital character, and a
trust-region step cap on the orbital rotation angle stabilises high-order
excited-state saddles (e.g.\ double excitations).
The text is organised as:
(\S\ref{sec:background}) mathematical background,
(\S\ref{sec:algorithm}) the full SOSCF workflow,
(\S\ref{sec:implementation}) the DFTB+ variables, subroutines, input
keywords, and output files that were added.
\end{abstract}

\tableofcontents
\bigskip
\hrule
\bigskip

%=====================================================================
\section{Introduction and motivation}

Ground-state SCF in DFTB+ is solved by self-consistent charge (SCC)
iteration with Broyden/DIIS charge mixing. For \emph{excited states}
at the mean-field level one wants a specific \emph{non-aufbau} orbital
occupation --- e.g.\ HOMO $\to$ LUMO --- to be the self-consistent
solution. At such a stationary point the orbital Hessian has one or
more negative eigenvalues (the target is a saddle, not a minimum), so
first-order SCF methods generically fail: simple mixing collapses into
the ground state, because descent along a negative-Hessian direction
\emph{is} the ground-state direction.

Two ideas solve this:
\begin{enumerate}[label=(\arabic*),leftmargin=*]
\item \textbf{Level-shifted Delta-SCF (STEP)} gives a warm start that
lives near the desired saddle by artificially shifting the
``to-be-virtual'' orbital eigenvalues so that the aufbau principle on
the shifted spectrum picks the target filling.
\item \textbf{Second-order SCF (SOSCF)}, a quasi-Newton method on the
manifold of orbital rotations, converges the excited state tightly by
taking Newton steps with a \emph{signed} inverse Hessian, which points
toward saddles by construction.
\end{enumerate}
Neither alone is sufficient: STEP is only contractive near its
artificial fixed point and cannot converge a saddle tightly, while
SOSCF diverges from a poor guess due to the initial-Hessian quality.
Combined with Initial MOM (IMOM) orbital tracking and a trust-region
step control, they form a robust Delta-SCF engine.

%=====================================================================
\section{Background and equations}\label{sec:background}

\subsection{SCC-DFTB ground-state SCF}

DFTB+ solves a generalised eigenvalue problem of Roothaan--Hall type,
iterated to self-consistency in the atomic charges. We summarise the
ingredients and the iteration structure here; \cite{Elstner1998,DFTBplus}
contain full derivations.

\paragraph{The SCC-DFTB Hamiltonian.}
The AO-basis Hamiltonian is split into a charge-independent part
$\mathbf{H}^{0}$ (from Slater--Koster tables, evaluated at neutral
reference densities) and a charge-dependent SCC correction:
\begin{equation}
H_{\mu\nu}[\qq] \;=\; H^{0}_{\mu\nu}
    \;+\; \tfrac12\, S_{\mu\nu}\, \sum_{\xi\in\text{atoms}}
     \big(\gamma_{\alpha(\mu)\,\xi} + \gamma_{\alpha(\nu)\,\xi}\big)\,
     \Delta q_\xi,
\label{eq:dftbH}
\end{equation}
where $\mu\in\alpha(\mu)$ labels the atom carrying AO $\mu$,
$\gamma_{\alpha\beta}$ is the DFTB gamma-function (a
distance-dependent Coulomb-like kernel regularised by the atomic
Hubbard parameters), and $\Delta q_\xi = q_\xi - q_\xi^{0}$ is the
deviation of the current Mulliken charge on atom $\xi$ from its
reference (neutral-atom) value. For spin-polarised calculations an
additional onsite spin coupling enters $\HH[\qq]$; it is treated
independently for each spin channel and does not alter the structure
below.

\paragraph{One-electron equation.}
At fixed $\qq$, the MOs are obtained from the generalised eigenvalue
problem
\begin{equation}
\FF[\qq]\,\CC \;=\; \SS\,\CC\,\bm{\varepsilon},\qquad
\CC^{\top}\SS\,\CC = \II,
\label{eq:dftbGEP}
\end{equation}
where $\SS$ is the AO overlap. Column $i$ of $\CC$ contains the MO
coefficients $C_{\mu i}$ of $\ket{\psi_i}=\sum_\mu C_{\mu i}\ket{\mu}$,
and $\varepsilon_i$ is the corresponding orbital energy. Equation
\eqref{eq:dftbGEP} is the DFTB analogue of the Roothaan--Hall equation
$\mathbf{F}\mathbf{C}=\mathbf{S}\mathbf{C}\bm{\varepsilon}$ of
Hartree--Fock theory, with $\mathbf{F}[\qq]$ replaced by the SCC-DFTB
Hamiltonian \eqref{eq:dftbH}.

\paragraph{Charge update and Mulliken analysis.}
At each SCC iteration the charges are recomputed from the MOs by a
Mulliken population analysis
\begin{equation}
q_\alpha \;=\; \sum_i f_i \!\! \sum_{\mu\in\alpha}\sum_\nu
       C_{\mu i}\, S_{\mu\nu}\, C_{\nu i}
\;=\; \sum_i f_i \!\! \sum_{\mu\in\alpha} (\mathbf{P}\SS)_{\mu\mu},
\label{eq:mulliken}
\end{equation}
with $\mathbf{P}_{\mu\nu}=\sum_i f_i C_{\mu i}C_{\nu i}$ the AO density
matrix and $f_i$ the occupation numbers (aufbau: lowest $n_\uparrow /
n_\downarrow$ per spin in the colinear unrestricted case, or lowest
$n_\mathrm{e}/2$ for closed shell). $\Delta\qq=\qq-\qq^{0}$ then feeds
back into Eq.~\eqref{eq:dftbH} for the next iteration.

\paragraph{SCF cycle (Roothaan--Hall style).}
One SCC iteration $k$ consists of:
\begin{enumerate}[leftmargin=*,nosep]
\item Build $\HH[\qq^{(k-1)}]$ from Eq.~\eqref{eq:dftbH}.
\item Solve Eq.~\eqref{eq:dftbGEP} $\to \CC^{(k)}, \bm{\varepsilon}^{(k)}$.
\item Apply aufbau to obtain $f_i^{(k)}$; build $\mathbf{P}^{(k)}$.
\item Compute $\qq^{(k,\mathrm{out})}$ from Eq.~\eqref{eq:mulliken}.
\item Convergence: $\mathrm{sccErrorQ}=\|\qq^{(k,\mathrm{out})}-\qq^{(k-1)}\|_\infty < \mathtt{SCCTolerance}$?
\item Else update the input charges by damped mixing
$\qq^{(k)} = \mathrm{Mixer}(\qq^{(k-1)},\,\qq^{(k,\mathrm{out})}-\qq^{(k-1)})$ (Broyden/DIIS); return to step~1.
\end{enumerate}
The fixed point $\qq^*$ of this map satisfies $\qq^*=\qq(\CC(\qq^*))$
and the MOs it produces diagonalise $\FF[\qq^*]$ in the $\SS$-metric.

\emph{All machinery in the remainder of this document operates on the
unrestricted case; each spin channel $\sigma\in\{\alpha,\beta\}$ is
treated independently, with its own $\FF^\sigma$, $\CC^\sigma$,
$f^\sigma$, and occ/virt partition.}

\subsection{Delta-SCF via State Targeted Energy Projection (STEP)}
\label{sec:step}

STEP is due to Carter-Fenk and Herbert~\cite{CarterFenkHerbert2020}.
It uses a level shift applied through the AO-basis virtual-density
operator $\mathbf{S}\mathbf{Q}\mathbf{S}$ to bias Aufbau selection
toward the target non-Aufbau configuration.

\paragraph{Non-aufbau target filling.}
Starting from the ground-state (aufbau) filling
$f^{(\mathrm{GS})}_i\in\{0,1\}$, the user specifies which orbital
indices to depopulate (``ExcitedFrom'') and which to populate
(``ExcitedTo'') to define the target filling $f^{(\mathrm{T})}_i$. For
example, for HOMO~$\to$~LUMO in a system with 4 occupied $\alpha$
orbitals, $f^{(\mathrm{T})}=\{1,1,1,0,1,0\}$.

\paragraph{Level-shift operator.}
Let $\{\ket{\psi^{(k)}_r}\}$ be the current MO set at iteration $k$,
with MO coefficients $\ket{\psi^{(k)}_r}=\sum_\mu C^{(k)}_{\mu r}\ket{\mu}$
in the AO basis $\{\ket{\mu}\}$, partitioned into the (current)
occupied set $\mathrm{occ}^{(k)}$ and virtual set $\mathrm{virt}^{(k)}$.
Throughout this section we write $a\in\mathrm{virt}^{(k)}$ for a
virtual-orbital index at iteration $k$, and correspondingly
$i\in\mathrm{occ}^{(k)}$; the occ/virt partition itself is implementation-specific
(see the code note below).

Define the virtual projector
\begin{equation}
\hat{Q}^{(k)} \;=\; \sum_{a\,\in\,\mathrm{virt}^{(k)}}
\ket{\psi^{(k)}_a}\bra{\psi^{(k)}_a},
\qquad
\hat{Q}^{(k)}\ket{\psi^{(k)}_r}\;=\;\begin{cases}
\ket{\psi^{(k)}_r}, & r\in\mathrm{virt}^{(k)},\\[2pt]
0, & r\in\mathrm{occ}^{(k)}.
\end{cases}
\end{equation}
The level-shifted Fock operator is
\begin{equation}
\hat{F}'^{(k)} \;=\; \hat{F}[\qq^{(k)}] \;+\; \eta\,\hat{Q}^{(k)},
\qquad \eta>0,
\label{eq:stepOp}
\end{equation}
so in the MO basis
\begin{equation}
\Braket{\psi^{(k)}_r|\hat{F}'^{(k)}|\psi^{(k)}_s}
\;=\; F^{(k)}_{rs} \;+\; \eta\,\delta_{rs}\,[\![\,r\!\in\!\mathrm{virt}^{(k)}\,]\!].
\end{equation}
Each currently-virtual eigenvalue is thus raised by $\eta$ while the
occupied eigenvalues and the occ--virt block of $\hat{F}$ in the MO
basis are unchanged.

For the generalised eigenvalue problem in the (non-orthogonal) AO
basis, we take AO matrix elements of Eq.~\eqref{eq:stepOp}. Using
$\Braket{\mu|\psi^{(k)}_a}=\sum_\nu S_{\mu\nu}C^{(k)}_{\nu a}=(\SS\CC^{(k)})_{\mu a}$,
\begin{equation}
\Braket{\mu|\hat{Q}^{(k)}|\nu}
\;=\; \sum_{a\,\in\,\mathrm{virt}^{(k)}}
\Braket{\mu|\psi^{(k)}_a}\Braket{\psi^{(k)}_a|\nu}
\;=\; \big[\SS\,\mathbf{Q}^{(k)}\,\SS\big]_{\mu\nu},
\qquad
Q^{(k)}_{\mu\nu}\;=\!\!\sum_{a\,\in\,\mathrm{virt}^{(k)}}\!
C^{(k)}_{\mu a}\,C^{(k)}_{\nu a}.
\end{equation}
The $\SS$-sandwich arises from AO non-orthogonality
($\Braket{\mu|\nu}=S_{\mu\nu}\neq\delta_{\mu\nu}$); dropping it would
be equivalent to $\bra\mu\hat Q\ket\nu=\mathbf Q_{\mu\nu}$, which only
holds on an orthonormal basis. The AO-basis representation of
Eq.~\eqref{eq:stepOp} is therefore
\begin{equation}
\boxed{\widehat{\FF}^{(k)} \;=\; \FF[\qq^{(k)}] \;+\; \eta\,\SS\,\mathbf{Q}^{(k)}\,\SS,}
\label{eq:step}
\end{equation}
which is the form solved in \cite[Eq.~6]{CarterFenkHerbert2020} and
implemented in \texttt{buildAndDiagDeltaScfRealHam}.

\paragraph{Code note: identifying $\mathrm{virt}^{(k)}$.}
The implementation does not store an integer list of virtual indices;
instead it inspects the real-valued filling array
$f^{(k)}_r\in[0,1]$ (integer $\{0,1\}$ for closed/open-shell
non-aufbau fillings here) and treats
\begin{equation}
r\in\mathrm{virt}^{(k)} \;\Longleftrightarrow\; f^{(k)}_r<\tfrac12.
\end{equation}
The threshold $\tfrac12$ is a robust tie-breaker between integer 0 and
1 fillings; it plays no mathematical role in the theory above. Inside
\texttt{buildAndDiagDeltaScfRealHam} this shows up as the loop
\texttt{if (targetFilling(iOrb,1,iSpin) < 0.5\_dp)} that selects the
virtual columns from \texttt{eigvecsReal} before forming
$\mathbf{Q}^{(k)}$.

Following the paper, we anchor $\eta$ to the HOMO--LUMO gap:
\begin{equation}
\eta \;=\; |\varepsilon_\mathrm{LUMO}^{(\mathrm{GS})}
           -\varepsilon_\mathrm{HOMO}^{(\mathrm{GS})}| \,+\, \varepsilon'
\label{eq:stepEta}
\end{equation}
with the default $\varepsilon'=0.1$~Ha. The gap-anchored value is
just large enough to force the desired occupation at every SCF
iteration without over-restricting occ--virt rotations, which a very
large $\eta$ would do. Per the paper, $\eta$ is determined \emph{once}
from the initial (non-Aufbau) orbitals and held fixed thereafter.
Alternatively a user-specified fixed $\eta$ (input keyword
\texttt{ShiftEnergy}) may be supplied.

\paragraph{STEP iteration.}
\begin{enumerate}
\item Build $\FF[\qq^{(k)}]$ and $\mathbf{Q}^{(k)}$ from the current
virtual columns of $\CC^{(k)}$, form
$\widehat{\FF}^{(k)}=\FF+\eta\,\SS\,\mathbf{Q}^{(k)}\,\SS$, and solve
\[
\widehat{\FF}^{(k)}\,\CC^{(k+1)} \;=\; \SS\,\CC^{(k+1)}\,\widehat{\bm{\varepsilon}}^{(k+1)}.
\]
\item Assign the $n_\mathrm{e}$ lowest shifted eigenvalues as occupied
(Aufbau on the shifted spectrum). This is the STEP prescription --- no
MOM/IMOM match is required in the original paper.
\item Restore physical eigenvalues by subtracting $\eta$ from the new
virtual indices:
$\varepsilon^{(k+1)}_a = \widehat{\varepsilon}^{(k+1)}_a - \eta$ for
$a\in\mathrm{virt}^{(k+1)}$.
\item Compute new density and new $\qq^{(k+1)}$; update charges
(direct substitution or Broyden).
\item Repeat until the \emph{physical} orbital gradient (not the
shifted one) satisfies $\max|g_{ia}|<\vartheta_\mathrm{SOSCF}$.
\end{enumerate}
Because the MO-basis representation of $\eta\,\SS\mathbf{Q}\SS$ is
diagonal on the (current) virtual block and zero elsewhere, the
occ--virt block of the MO-transformed Fock matrix is unchanged by the
shift. Consequently the \emph{physical} orbital gradient at the
shifted fixed point equals the gradient for $\FF$ alone: the shifted
fixed point coincides with a stationary point of the \emph{physical}
energy in the occ/virt partition set by $f^{(\mathrm{T})}$.

\subsection{Initial Maximum Overlap Method (IMOM)}
\label{sec:imom}

The original STEP paper~\cite{CarterFenkHerbert2020} uses Aufbau on
the shifted spectrum alone, without MOM/IMOM, and positions STEP as
an alternative to MOM-style overlap tracking. In practice, however,
for non-adjacent or multi-index excitations the Aufbau assignment on
$\widehat{\bm{\varepsilon}}^{(k+1)}$ can become ambiguous when target
virtual levels drift through near-degeneracy with occupied levels, so
we offer IMOM as an optional overlap-based reassignment layered on top
of STEP (\texttt{UseIMOM}, default \texttt{Yes}).

Initial MOM fixes a \emph{frozen} reference set of eigenvectors
$\CC^{(\mathrm{ref})}$ (the ground-state eigenvectors at loop entry)
together with a fixed reference filling $f^{(\mathrm{ref})}$ identifying
the reference virtual set $\mathrm{virt}^{(\mathrm{ref})}$ (in the code: orbitals with $f^{(\mathrm{ref})}<\tfrac12$).
At each iteration, given the new eigenvectors $\CC^{(k+1)}$, the
overlap matrix
\begin{equation}
M_{ij} \;=\; \langle \psi^{(k+1)}_i \mid \widehat{S} \mid \psi^{(\mathrm{ref})}_j \rangle
\;=\;(\CC^{(k+1)})^{\top} \SS\, \CC^{(\mathrm{ref})}
\end{equation}
is computed, and filling is re-assigned by \textbf{greedy exclusive
maximum-overlap} matching: for each $j \in \mathrm{virt}^{(\mathrm{ref})}$, the new
orbital index
\begin{equation}
i^\ast(j) \;=\; \operatorname*{argmax}_{i \notin \mathrm{assigned}}\; |M_{ij}|^2
\end{equation}
is marked virtual; all remaining new indices are marked occupied.
This reassignment replaces aufbau in the STEP step~(2) above and is
carried through the entire SOSCF iteration as the filling rule.

\subsection{Second-Order SCF (SOSCF)}
\label{sec:soscf-theory}

\paragraph{Exponential parameterisation.}
Write the new MOs as a unitary rotation of the current ones:
\begin{equation}
\CC^\mathrm{new} \;=\; \CC^\mathrm{old}\, \UU, \qquad
\UU \;=\; \exp(\kkappa),
\end{equation}
with $\kkappa$ a real antisymmetric matrix ($\kkappa^\top=-\kkappa$).
The only independent variables are the occ--virt block
$x_{ia}=\kappa_{ia}$ (with $i$ in the occupied and $a$ in the virtual
set of the current filling), because rotations within the occupied or
virtual block do not change the density matrix
$\rho=\sum_i f_i|\psi_i\rangle\langle\psi_i|$.

For one spin channel, stacking $x_{ia}$ into a vector $\xx$ of length
$n_\mathrm{oo} \equiv n_\mathrm{occ}\cdot n_\mathrm{virt}$, the energy
is a function $E(\xx)$ and we optimise it on
$\xx\in\mathbb{R}^{n_\mathrm{oo}}$ with
\begin{equation}
\kappa_{ia}\!=\!x_{ia},\quad \kappa_{ai}\!=\!-x_{ia},\quad
\kappa_{ij}\!=\!\kappa_{ab}\!=\!0.
\end{equation}

\paragraph{Orbital gradient.}
At $\xx=0$ the gradient reads
\begin{equation}
g_{ia} \;=\; \frac{\partial E}{\partial x_{ia}}\bigg|_{\xx=0}
\;=\; 4\,[\,\CC^\top \FF \,\CC\,]_{ia}
\;\equiv\; 4\, F^\mathrm{MO}_{ia}.
\label{eq:grad}
\end{equation}
We derive this in four steps: (i) write the SCC-DFTB energy as a
functional of the density matrix, (ii) propagate the orbital rotation
into the density, (iii) differentiate with respect to $\kappa_{rs}$,
(iv) apply the chain rule from $\kappa$ to $x$ and collect the factor.

\emph{(i) Energy as a functional of $\mathbf{P}$.}
For the single Slater determinant defined by filling $f$ and MO matrix
$\CC$, the AO density matrix is
\begin{equation}
P_{\mu\nu} \;=\; \sum_p f_p\, C_{\mu p}\, C_{\nu p}.
\label{eq:densitymat}
\end{equation}
The SCC-DFTB total energy \eqref{eq:dftbH} is a functional
$E[\mathbf{P}]$ whose gradient with respect to $\mathbf{P}$ \emph{is},
by construction of the Fock/Kohn--Sham matrix, the Hamiltonian itself:
\begin{equation}
\frac{\partial E}{\partial P_{\mu\nu}}\bigg|_{\text{const }\qq\text{-dep}} \;=\; H_{\mu\nu}[\qq],
\qquad \text{with } \qq \text{ also a functional of } \mathbf{P}.
\label{eq:dEdP}
\end{equation}
The crucial simplification for SCF-type energies is that, precisely
because $\FF$ is built to be $\partial E/\partial\mathbf{P}$, the
derivatives of $\qq$ with respect to $\mathbf{P}$ cancel out in the
full $\mathrm{d}E/\mathrm{d}\mathbf{P}$ (this is the statement of the
variational property / generalised Brillouin theorem). We may
therefore treat $\FF$ as constant under orbital rotation when computing
the first derivative of $E$, a standard property of Hartree--Fock-- and
DFT-type energies \cite[Ch.~10]{HelgakerBook}.

\emph{(ii) Rotation of $\mathbf{P}$.}
Under $\CC\to\CC\,\UU$ with $\UU=\exp(\kkappa)$, the density matrix
for fixed $f$ transforms as
$\mathbf{P}\to\CC\,\UU\,\mathbf{f}\,\UU^{\top}\CC^{\top}$ where
$\mathbf{f}=\mathrm{diag}(f_p)$. Expanding $\UU=\II+\kkappa+O(\kkappa^2)$,
\begin{equation}
P^\mathrm{MO}_{pq}(\kkappa)
\;\equiv\; (\UU\,\mathbf{f}\,\UU^{\top})_{pq}
\;=\; f_p\delta_{pq} + (f_q - f_p)\kappa_{pq} + O(\kkappa^2).
\label{eq:Pmo-rot}
\end{equation}
The MO-basis density change is thus
$\Delta P^\mathrm{MO}_{pq} = (f_q-f_p)\kappa_{pq}$ to first order.

\emph{(iii) Differentiation w.r.t.\ $\kappa_{rs}$.}
Writing $E$ in the MO basis, $E = \sum_{pq} P^\mathrm{MO}_{pq} F^\mathrm{MO}_{qp}/2 + \text{(terms independent of } \kkappa)$ for the SCF part and using
\eqref{eq:dEdP}, and holding $F^\mathrm{MO}$ fixed (the variational
property above),
\begin{equation}
\frac{\partial E}{\partial \kappa_{rs}}\bigg|_{\kkappa=0}
\;=\; \sum_{pq}\frac{\partial P^\mathrm{MO}_{pq}}{\partial \kappa_{rs}}\,F^\mathrm{MO}_{qp}
\;=\; (f_s - f_r)\,F^\mathrm{MO}_{rs} \;+\; (f_r - f_s)\,F^\mathrm{MO}_{sr}
\;=\; 2(f_s-f_r)\,F^\mathrm{MO}_{rs},
\label{eq:dEdkappa}
\end{equation}
where the last equality uses $F^\mathrm{MO}_{sr}=F^\mathrm{MO}_{rs}$
(Hermiticity with real MOs). For $r=i\in\mathrm{occ}$, $s=a\in\mathrm{virt}$
with integer fillings $f_i=1,\;f_a=0$ this gives
$\partial E/\partial\kappa_{ia}\big|_0 = -2\,F^\mathrm{MO}_{ia}$, and
$\partial E/\partial\kappa_{ai}\big|_0 = +2\,F^\mathrm{MO}_{ai} = +2\,F^\mathrm{MO}_{ia}$.

\emph{(iv) Chain rule $\kappa\to x$.}
Since the independent variable is $x_{ia}$ with
$\kappa_{ia}=x_{ia}$ and $\kappa_{ai}=-x_{ia}$,
\begin{equation}
\frac{\partial E}{\partial x_{ia}}\bigg|_{0}
\;=\; \frac{\partial E}{\partial \kappa_{ia}}\bigg|_{0}\frac{\partial \kappa_{ia}}{\partial x_{ia}}
\;+\; \frac{\partial E}{\partial \kappa_{ai}}\bigg|_{0}\frac{\partial \kappa_{ai}}{\partial x_{ia}}
\;=\; (-2\,F^\mathrm{MO}_{ia})(+1) + (+2\,F^\mathrm{MO}_{ia})(-1)
\;=\; -4\,F^\mathrm{MO}_{ia}.
\end{equation}
The sign depends on the sign convention for $\kappa$ in the
exponential; the code uses the opposite convention $\CC_\mathrm{new}=\CC\,\UU(-\kkappa)$-equivalent sign
such that $g_{ia}=+4\,F^\mathrm{MO}_{ia}$ as in Eq.~\eqref{eq:grad} (see
\texttt{soscf.F90:283}). The magnitude~4 is a factor-of-two product
from antisymmetry of $\kkappa$ (both the $ia$ and $ai$ entries
contribute) times the factor $2$ in the density-matrix transformation
\eqref{eq:Pmo-rot}, and is kept consistent with the factor~2 in the
initial diagonal Hessian \eqref{eq:hess-diag} so that the quasi-Newton
step $\mathrm{D}\xx=-\HH^{-1}\gg$ has the standard
$F^\mathrm{MO}_{ia}/(\varepsilon_a-\varepsilon_i)$ scale. The
stationary point $g_{ia}=0 \Leftrightarrow F^\mathrm{MO}_{ia}=0$
(block-diagonal MO Fock) is independent of this overall prefactor.

\paragraph{Physical interpretation.}
$g_{ia}\propto F^\mathrm{MO}_{ia}$ is the off-diagonal occ--virt
matrix element of the Fock/Kohn--Sham operator in the current MO
basis. At an SCF stationary point this block is exactly zero, so
$g_{ia}=0$ means the current MOs already diagonalise $\FF$ in the
occ/virt partition chosen by $f^{(\mathrm{T})}$. Driving
$\max_{ia}|g_{ia}|\to 0$ is therefore equivalent to satisfying
Eq.~\eqref{eq:dftbGEP} on the non-Aufbau manifold.

\paragraph{Orbital Hessian (diagonal approximation).}
The diagonal element of the orbital Hessian in the ORCA
convention is
\begin{equation}
H_{ia,ia}
\;=\; 2\,(f_a - f_i)(\varepsilon_i - \varepsilon_a)
\;=\; 2\,(\varepsilon_a - \varepsilon_i) \quad \text{for integer filling}
\label{eq:hess-diag}
\end{equation}
where the final equality uses $f_\mathrm{occ}=1,\,f_\mathrm{virt}=0$.
The corresponding initial inverse Hessian (used at the first SOSCF
step) is
\begin{equation}
\big[H^{-1}\big]_{ia,ia}^{(0)}
\;=\;\frac{\operatorname{sign}(\Delta_{ia})}
{\max\!\big(2\,|\Delta_{ia}|,\; 10^{-3}\big)},
\qquad \Delta_{ia}\!=\!\varepsilon_a-\varepsilon_i.
\label{eq:invH-init}
\end{equation}
The \emph{signed} form \eqref{eq:invH-init} is essential for
excited-state saddles: when the target has a negative-curvature mode
(orbital $a$ virtual in target but $\varepsilon_a<\varepsilon_i$ in
the ground state), $\Delta_{ia}<0$ and the inverse-Hessian element
becomes negative, so that the Newton step
$\mathrm{D}x=-H^{-1}g$ points \emph{uphill} along that
direction --- exactly what is needed to converge to a saddle.

\paragraph{Quasi-Newton step.}
At iteration $n$,
\begin{equation}
\mathrm{D}\xx_n \;=\; -\,\HH_n^{-1}\,\gg_n,
\label{eq:newton-step}
\end{equation}
followed by a trust-region clip (\S\ref{sec:trust-region} below) and
orbital rotation
\begin{equation}
\CC_{n+1} \;=\; \CC_n\,\exp(\kkappa_n), \qquad
\kkappa_n \leftarrow \text{stack of } \mathrm{D}x_{ia}.
\end{equation}

\paragraph{Cayley approximant for $\exp(\kkappa)$.}
The unitary factor is computed via the Pad\'e[1,1] (Cayley) form
\begin{equation}
\UU \;=\; (\II-\tfrac12\kkappa)^{-1}(\II+\tfrac12\kkappa),
\label{eq:cayley}
\end{equation}
which is \emph{exactly} orthogonal for any real antisymmetric
$\kkappa$ (no truncation error in orthogonality), is second-order
accurate in $\kkappa$, and reduces to one linear solve per iteration.
This avoids the Taylor truncation error and orthogonality drift of
$\UU \approx \II + \kkappa + \tfrac12\kkappa^2 + \dots$.

\subsection{L-SR1 inverse-Hessian update}
\label{sec:lsr1}

The true orbital Hessian is $O(n_\mathrm{oo}^2)$ to store and expensive
to evaluate, so SOSCF maintains an approximation that is updated from
secant information after every iteration. The choice of update rule is
\emph{not} cosmetic for excited-state $\Delta$-SCF: it is dictated by
the sign structure of the Hessian at the target stationary point.

\paragraph{Why not BFGS: the indefinite-Hessian obstruction.}
Standard quasi-Newton optimisation methods like BFGS are designed for
\emph{minimisation} of convex objectives and are built to enforce
positive definiteness of the approximate Hessian. Concretely, the
BFGS update of the inverse Hessian,
\begin{equation}
\HH^{-1,\mathrm{BFGS}}_{n+1}
\;=\; \Big(\II - \frac{\bm{\delta}_n\bm{\gamma}_n^{\top}}{\bm{\gamma}_n^{\top}\bm{\delta}_n}\Big)
      \HH^{-1}_n
      \Big(\II - \frac{\bm{\gamma}_n\bm{\delta}_n^{\top}}{\bm{\gamma}_n^{\top}\bm{\delta}_n}\Big)
 \;+\; \frac{\bm{\delta}_n\bm{\delta}_n^{\top}}{\bm{\gamma}_n^{\top}\bm{\delta}_n},
\label{eq:bfgs}
\end{equation}
preserves $\HH^{-1}_n\succ 0\Rightarrow\HH^{-1}_{n+1}\succ 0$ whenever
the curvature condition $\bm{\gamma}_n^{\top}\bm{\delta}_n>0$ is
satisfied (and damping is used when it is not). Applied to an
excited-state saddle this is catastrophic: at the target, the true
orbital Hessian has at least one \emph{negative} eigenvalue along the
$\mathrm{GS}\to\mathrm{ES}$ de-excitation direction, and the
saddle-seeking Newton step $\mathrm{D}\xx=-\HH^{-1}\gg$ requires
$\HH^{-1}$ to inherit that negative sign. A BFGS-maintained
$\HH^{-1}\succ 0$ always yields descent directions, so
$\mathrm{D}\xx=-\HH^{-1}\gg$ points \emph{downhill} along the GS
direction --- i.e., toward variational collapse to the Aufbau ground
state. The same objection rules out L-BFGS, damped-BFGS, and any
other positive-definiteness-enforcing variant.

\paragraph{Why SR1: signature-preserving rank-1 update.}
The symmetric rank-$1$ (SR1) update is the unique symmetric update of
rank~1 that satisfies the secant equation, and it \emph{does not}
enforce any sign on the curvature it generates. If the secant pair
$(\bm{\delta}_n,\bm{\gamma}_n)$ reveals negative curvature
($\bm{j}_n^{\top}\bm{\gamma}_n<0$, with $\bm{j}_n$ defined below), the
SR1 correction subtracts rather than adds to $\HH^{-1}$, so the
approximate inverse Hessian can (and does) become indefinite whenever
the true Hessian is. This is the standard reason SR1 is recommended
over BFGS for trust-region methods that must converge to saddle
points \cite[Ch.~6]{NocedalWright}.

Using SR1 with a \emph{signed} initial diagonal inverse Hessian
(Eq.~\eqref{eq:invH-init}) gives SR1 the correct starting signature
in each occ--virt direction; the rank-1 updates then refine, but do
not wash out, the negative-curvature modes along which the target is
a saddle.

\paragraph{Secant pair and secant condition.}
After the step $\mathrm{D}\xx_n$ and gradient recomputation,
\begin{equation}
\bm{\delta}_n \;\equiv\; \mathrm{D}\xx_n, \qquad
\bm{\gamma}_n \;\equiv\; \gg_{n+1} - \gg_n.
\end{equation}
The standard quasi-Newton \emph{secant condition} on the inverse
Hessian is that $\HH^{-1}_{n+1}$ should map the observed gradient
change back to the observed step:
\begin{equation}
\HH^{-1}_{n+1}\,\bm{\gamma}_n \;=\; \bm{\delta}_n.
\label{eq:secant}
\end{equation}
This is the first-order Taylor statement
$\bm{\gamma}_n = \HH_{n+1}\bm{\delta}_n+O(\|\bm\delta\|^2)$ with $\HH = (\HH^{-1})^{-1}$, and any quasi-Newton update of $\HH^{-1}$
must satisfy it.

\paragraph{Derivation of the SR1 update (inverse-Hessian form).}
We seek the simplest update of $\HH^{-1}$ that (a) is symmetric,
(b) is rank~1, (c) satisfies \eqref{eq:secant}, and (d) does not
constrain the sign of the curvature generated. The rank-1 symmetric
ansatz is
\begin{equation}
\HH^{-1}_{n+1} \;=\; \HH^{-1}_n \;+\; \sigma\,\mathbf{v}\,\mathbf{v}^{\top},
\qquad \sigma\in\{+1,-1\},\ \mathbf{v}\in\mathbb{R}^{n_\mathrm{oo}}.
\label{eq:sr1ansatz}
\end{equation}
Applying the secant condition \eqref{eq:secant} gives
\begin{align}
\HH^{-1}_n\,\bm{\gamma}_n \;+\; \sigma\,(\mathbf{v}^{\top}\bm{\gamma}_n)\,\mathbf{v}
\;&=\;\bm{\delta}_n \notag\\
\Longrightarrow\quad
\sigma\,(\mathbf{v}^{\top}\bm{\gamma}_n)\,\mathbf{v}
\;&=\;\bm{\delta}_n \;-\; \HH^{-1}_n\,\bm{\gamma}_n
\;\equiv\;\bm{j}_n.
\label{eq:sr1-v-parallel-j}
\end{align}
Equation~\eqref{eq:sr1-v-parallel-j} says $\mathbf{v}$ is parallel to
$\bm{j}_n$, i.e.\ $\mathbf{v}=\lambda\,\bm{j}_n$ for some scalar
$\lambda$. Substituting,
\begin{equation}
\sigma\,\lambda^2\,(\bm{j}_n^{\top}\bm{\gamma}_n)\,\bm{j}_n
\;=\; \bm{j}_n
\;\Longrightarrow\;
\sigma\,\lambda^2\;=\;\frac{1}{\bm{j}_n^{\top}\bm{\gamma}_n}.
\end{equation}
The combined scalar $\sigma\lambda^2$ is uniquely determined; absorbing
it gives the SR1 inverse-Hessian update
\begin{equation}
\boxed{\HH^{-1}_{n+1} \;=\; \HH^{-1}_n \;+\;
\frac{\bm{j}_n\,\bm{j}_n^{\top}}{\bm{j}_n^{\top}\bm{\gamma}_n},
\qquad \bm{j}_n \;=\; \bm{\delta}_n - \HH^{-1}_n\,\bm{\gamma}_n.}
\label{eq:lsr1}
\end{equation}
The sign of the rank-1 correction is carried by the scalar
$\bm{j}_n^{\top}\bm{\gamma}_n$, which is \emph{not} constrained to be
positive: when the secant pair samples a negative-curvature direction
(the true Hessian has negative eigenvalues there), the denominator is
negative and $\HH^{-1}$ picks up a negative rank-1 correction along
$\bm{j}_n$. This is exactly the behaviour required to track a
saddle-shaped objective, and is the main technical reason SR1 (and
not BFGS) is used here.

\paragraph{Safeguards.}
Two numerical safeguards are applied in the code
(\texttt{updateInvHessianLSR1}):
\begin{itemize}[leftmargin=*,nosep]
\item \textbf{Denominator guard:} skip the update if
$|\bm{j}_n^{\top}\bm{\gamma}_n| < \tau\,\|\bm{j}_n\|\,\|\bm{\gamma}_n\|$
with $\tau=10^{-8}$. Equivalently, the rank-1 coefficient
$1/(\bm{j}_n^{\top}\bm{\gamma}_n)$ is rejected when $\bm{j}_n$ is
nearly orthogonal to $\bm{\gamma}_n$, which would otherwise divide by a
catastrophically small denominator and inflate $\|\HH^{-1}\|$.
This is the classic SR1 failure mode \cite[\S6.2]{NocedalWright}.
\item \textbf{Step-norm guard:} skip if $\|\bm{\delta}_n\| < 10^{-5}$.
A negligibly small step samples no curvature; the resulting
$(\bm{\delta},\bm{\gamma})$ is dominated by floating-point noise, and
the rank-1 update amplifies that noise as $\|\bm{\delta}\|^2/|\bm{j}^{\top}\bm{\gamma}|$.
\end{itemize}
When either guard fires, $\HH^{-1}_{n+1}=\HH^{-1}_n$ is used; the
missed update simply defers curvature refinement to the next
iteration.

\subsection{Trust-region step control}
\label{sec:trust-region}

The diagonal initial inverse Hessian \eqref{eq:invH-init} ignores the
off-diagonal Coulomb/exchange couplings between occ--virt pairs. For
high-order excited-state saddles (e.g.\ double excitations) the
un-scaled Newton step \eqref{eq:newton-step} can exceed $1$~rad,
invalidating the Cayley linearisation, scrambling IMOM assignments,
and poisoning subsequent L-SR1 updates.

We cap the per-iteration rotation magnitude:
\begin{equation}
\boxed{
\mathrm{D}\xx_n \leftarrow \mathrm{D}\xx_n \cdot
\min\!\left(1,\; \frac{\kappa_\mathrm{max}}{\max_{ia}|\mathrm{D}x_{ia}|}\right).
}
\label{eq:trust-cap}
\end{equation}
The step \emph{direction} is preserved, only the magnitude is
shrunk. Below the cap (e.g.\ at convergence) the rescaling is the
identity, so we recover full superlinear Newton convergence. The
default $\kappa_\mathrm{max}=0.5$~rad
($\sim 29^\circ$) is a standard trust-region setting for orbital
optimisation~\cite{Bacskay1981,HelgakerBook}. See
\S\ref{sec:results-dblexc} for the rationale of this value.

%=====================================================================
\section{Full algorithm}\label{sec:algorithm}

After the ground-state SCC loop, the driver dispatches on \texttt{tSoscf}:
\begin{itemize}[leftmargin=*,nosep]
\item If \texttt{tSoscf = .true.}: run \texttt{runPreSoscfDeltaScfLoop}
  (\S\ref{sec:algo-preSOSCF}, warm-start) then
  \texttt{runSoscfLoop} (\S\ref{sec:algo-soscf}, tight convergence).
\item If \texttt{tSoscf = .false.}: run \texttt{runDeltaScfStandaloneLoop},
  a level-shifted Delta-SCF that converges on the ground-state SCC
  criterion $\mathrm{sccErrorQ}<\mathtt{SCCTolerance}$ and returns that
  result directly (no orbital-rotation phase).
\end{itemize}

\subsection{Overall flow}

\begin{algorithm}[H]
\caption{DFTB+ excited-state Delta-SCF / SOSCF driver (see \texttt{main.F90:1371}).}
\label{alg:overall}
\begin{algorithmic}[1]
\State Run ground-state SCC loop to convergence (sccErrorQ $<$ SCCTolerance).
\State Save GS eigenvectors $\CC^{(\mathrm{GS})}$ and eigenvalues $\bm{\varepsilon}^{(\mathrm{GS})}$.
\State Build target filling $f^{(\mathrm{T})}$ from ExcitedFrom / ExcitedTo indices (\texttt{buildDeltaScfTargetFilling}).
\If{\texttt{tDeltaScf} \textbf{and} \texttt{tRealHS}}
  \If{\texttt{tSoscf}}
    \State \textbf{runPreSoscfDeltaScfLoop} $\to$ warm-start; exits on $\max|g|<\vartheta_\mathrm{SOSCF}$ (Alg.~\ref{alg:deltaScf}).
    \State \textbf{runSoscfLoop} $\to$ tight SOSCF; one Newton rotation per outer iter (Alg.~\ref{alg:soscf}). Refreshes $\bm{\varepsilon}\leftarrow\mathrm{diag}(\CC^{\top}\FF\CC)$ on exit.
  \Else
    \State \textbf{runDeltaScfStandaloneLoop} $\to$ exits on $\mathrm{sccErrorQ}<\mathtt{SCCTolerance}$.
  \EndIf
\EndIf
\State Write band.out (sorted ascending by $\varepsilon$ after SOSCF; unsorted otherwise) / detailed.out / charges.bin.
\end{algorithmic}
\end{algorithm}

\subsection{Pre-SOSCF: level-shifted Delta-SCF (STEP\,$\pm$\,IMOM)}
\label{sec:algo-preSOSCF}

\begin{algorithm}[H]
\caption{\texttt{runPreSoscfDeltaScfLoop} (main.F90:7733). \textbf{Exit
criterion:} $\max|g|<\vartheta_\mathrm{SOSCF}$ only; $\mathrm{sccErrorQ}$
is computed and printed for diagnostics but is \emph{not} an exit
condition (charge self-consistency is handled by the SOSCF outer loop).}
\label{alg:deltaScf}
\begin{algorithmic}[1]
\If{\texttt{UseIMOM}}
  \State $\CC^{(\mathrm{ref})} \gets \CC^{(\mathrm{GS})}$;
         $f^{(\mathrm{ref})}\gets f^{(\mathrm{T})}$
         \Comment{frozen IMOM reference}
\EndIf
\State Determine level shift $\eta$ \emph{once}: user-supplied
       \texttt{ShiftEnergy}, or
       $\eta = \varepsilon_\mathrm{LUMO}^{(\mathrm{GS})}
             - \varepsilon_\mathrm{HOMO}^{(\mathrm{GS})}
             + \varepsilon'$ with
       $\varepsilon'\equiv\mathtt{ShiftMargin}$
       (Eq.~\eqref{eq:stepEta}).
\State Initialise $f^{(0)} \gets f^{(\mathrm{T})}$
\For{$k = 1, 2, \dots, \mathrm{MaxSCC}$}
  \State {[a]} Build $\FF[\qq^{(k-1)}]$ (\texttt{getSccHamiltonian}).
  \State {[b,c]} Build $\mathbf{Q}^{(k-1)}=\sum_{a\in\mathrm{virt}^{(k-1)}} \CC^{(k-1)}_{:,a}(\CC^{(k-1)}_{:,a})^{\top}$ from current virtual columns and solve the level-shifted problem
         $(\FF + \eta\,\SS\mathbf{Q}^{(k-1)}\SS)\CC^{(k)}
         = \SS\CC^{(k)}\widehat{\bm{\varepsilon}}^{(k)}$ (\texttt{buildAndDiagDeltaScfRealHam}, Eq.~\eqref{eq:step}).
  \State {[d]} Aufbau on $\widehat{\bm{\varepsilon}}^{(k)}$ sets
         tentative $f^{(k)}$ (\texttt{getFillingsAndBandEnergies}).
  \If{\texttt{UseIMOM}}
    \State {[d']} compute $M=(\CC^{(k)})^\top \SS\, \CC^{(\mathrm{ref})}$; greedy virtual-index reassignment overrides $f^{(k)}$ (\texttt{getIMOMTargetFilling}).
  \EndIf
  \State {[e]} Restore physical eigenvalues: $\varepsilon^{(k)}_a \gets \widehat{\varepsilon}^{(k)}_a - \eta$ for $a\in\mathrm{virt}^{(k)}$ (\texttt{restoreDeltaScfEigenvalues}).
  \State {[e2]} Compute physical gradient $g^{(k)}_{ia} = 4\,[\,(\CC^{(k)})^\top\FF\CC^{(k)}\,]_{ia}$ from the \emph{unshifted} $\FF$ (\texttt{computeOrbGradient}).
  \If{$\max_{ia}|g^{(k)}_{ia}| < \vartheta_\mathrm{SOSCF}$}
     \State \textbf{exit (handoff to SOSCF)}
  \EndIf
  \State {[f--k]} Build density $\rho$, output charges $\qq^{(k,\mathrm{out})}$, energies.
  \State \hskip 1.0em Compute $\mathrm{sccErrorQ}=\|\qq^{(k,\mathrm{out})}-\qq^{(k-1)}\|_\infty$ (diagnostic only).
  \State {[l]} Charge update ($\qq^{(k)}\!:=\!\qq^\mathrm{in}$ for iter $k{+}1$):
     \textbf{if \texttt{UseMixer}}: Broyden/DIIS step on $\qq^{(k,\mathrm{out})}-\qq^{(k-1)}$;
     \textbf{else}: $\qq^{(k)}\gets\qq^{(k,\mathrm{out})}$.
\EndFor
\end{algorithmic}
\end{algorithm}

The companion routine \texttt{runDeltaScfStandaloneLoop} (used when
\texttt{tSoscf = .false.}) shares steps [a]--[l] but replaces the
$\max|g|$-based exit with the standard SCC criterion
$\mathrm{sccErrorQ}<\mathtt{SCCTolerance}$; no gradient is computed
there.

\subsection{SOSCF outer loop}
\label{sec:algo-soscf}

\begin{algorithm}[H]
\caption{\texttt{runSoscfLoop} (main.F90:8339). One Broyden/DIIS charge
update \emph{and} one quasi-Newton orbital rotation per outer
iteration; no inner diagonalisation.}
\label{alg:soscf}
\begin{algorithmic}[1]
\State \texttt{soscf\_init}: for each spin, partition occ/virt via
       $f^{(\mathrm{T})}$, allocate
       $\HH^{-1}\in\mathbb{R}^{n_\mathrm{oo}\times n_\mathrm{oo}}$, set
       $\HH^{-1} \gets \mathrm{diag}\big(\operatorname{sign}(\Delta_{ia})/\max(2|\Delta_{ia}|,10^{-3})\big)$
       per Eq.~\eqref{eq:invH-init}; record $\kappa_\mathrm{max}$.
\State Initialise orbital-tracking reference from the post-pre-SOSCF state:
       $\CC^{(\mathrm{ref})}\gets\CC$, $f^{(\mathrm{ref})}\gets f^{(\mathrm{T})}$.
       \Comment{\texttt{UseIMOM=.true.}: frozen; \texttt{.false.}: advanced at end of Step A each iter.}
\State Reset the Broyden/DIIS mixer history once (only if \texttt{UseMixer}).
\For{$n = 1, 2, \dots, \mathrm{MaxSCC}$}
  \State \textbf{Step A --- Orbital rotation at the current $\qq^\mathrm{in}$}:
  \State \qquad Rebuild $\FF\gets\FF[\qq^\mathrm{in}]$ from the current input charges (inherited from pre-SOSCF at $n=1$, otherwise from Step~B of iter $n{-}1$).
  \State \qquad \textbf{for each spin:}
  \State \qquad\qquad Compute gradient $\gg_{ia}=4[\,\CC^{\top}\FF\,\CC\,]_{ia}$ on occ--virt (\texttt{computeOrbGradient}); record $\max|g|$ for outer convergence.
  \State \qquad\qquad \textbf{if $n>1$}: L-SR1 update~\eqref{eq:lsr1} with $\bm{\delta}=\mathrm{D}\xx_{n-1}$, $\bm{\gamma}=\gg-\gg_{\mathrm{prev}}$ (\texttt{updateInvHessianLSR1}).
  \State \qquad\qquad Store $\gg_\mathrm{prev}\gets\gg$. Compute Newton step $\mathrm{D}\xx=-\HH^{-1}\gg$ with trust-region clip~\eqref{eq:trust-cap} (\texttt{getSoscfStep}).
  \State \qquad\qquad Store $\mathrm{D}\xx_\mathrm{prev}\gets\mathrm{D}\xx$. Build $\kkappa$ from $\mathrm{D}\xx$ (\texttt{buildKappa}), form Cayley $\UU$~\eqref{eq:cayley} (\texttt{computeCayleyExp}), rotate $\CC\gets\CC\,\UU$.
  \State \qquad\qquad Reassign $f^{(\mathrm{T})}$ by max-overlap of $\CC^\mathrm{new}$ against $(\CC^{(\mathrm{ref})},f^{(\mathrm{ref})})$ (\texttt{getIMOMTargetFilling}).
  \State \qquad \textbf{if \texttt{not UseIMOM}}: advance tracking reference
         $\CC^{(\mathrm{ref})}\gets\CC^\mathrm{new}$, $f^{(\mathrm{ref})}\gets f^{(\mathrm{T})}$ \Comment{classical MOM}
  \State \textbf{Step B --- Density, charges and energies at the post-rotation $\CC$}:
  \State \qquad Build $\rho$ from $\CC^\mathrm{new}$ and $f^{(\mathrm{T})}$; compute output charges $\qq^\mathrm{out}$ (Mulliken).
  \State \qquad Evaluate all energy contributions and $E_\mathrm{elec}$ at the rotated state (\texttt{calcEnergies}, \texttt{sumEnergies}).
  \State \qquad $\mathrm{sccErrorQ}\gets\|\qq^\mathrm{out}-\qq^\mathrm{in}\|_\infty$ (residual at $\CC^\mathrm{new}$ against the $\qq^\mathrm{in}$ used by Step~A of this iter).
  \State \qquad \textbf{if \texttt{UseMixer}}:
         $\qq^\mathrm{in}\gets\mathrm{Broyden}(\qq^\mathrm{in},\qq^\mathrm{out}-\qq^\mathrm{in})$;
         \textbf{else}: $\qq^\mathrm{in}\gets\qq^\mathrm{out}$. \Comment{advances $\qq^\mathrm{in}$ for Step~A of iter $n{+}1$}
  \State \textbf{Step C --- Convergence check}:
  \If{$\mathrm{sccErrorQ}<\mathtt{OuterSCCTolerance}$
      \textbf{and} $\max|\gg|<\vartheta_\mathrm{SOSCF}$}
     \State \textbf{exit}
  \EndIf
\EndFor
\State Refresh $\bm{\varepsilon}\leftarrow\mathrm{diag}(\CC^{\top}\FF\CC)$ so downstream output (band.out, detailed.out) is consistent with the rotated MOs.
\end{algorithmic}
\end{algorithm}

\paragraph{Post-processing.}
After the SOSCF outer loop exits, the eigenvalue array
$\bm{\varepsilon}$ is stale (SOSCF rotates the MOs but does not
re-diagonalise). We refresh it by
$\varepsilon_i \leftarrow \langle\CC_i|\FF|\CC_i\rangle$
(the diagonal of $F^\mathrm{MO}$), which at the SOSCF stationary
point --- where $F^\mathrm{MO}$ is block-diagonal between occ and
virt --- is the physical orbital energy up to a within-block rotation.

%=====================================================================
\section{Implementation in DFTB+}\label{sec:implementation}

\subsection{Input keywords (HSD)}

All excited-state keywords live inside the
\texttt{Hamiltonian=DFTB\{\dots\}} block as two sub-blocks:

\begin{lstlisting}[style=hsd,caption={Input example: formaldehyde HOMO-1 $\to$ LUMO+1 / HOMO $\to$ LUMO+2 double excitation (illustrates overrides; keywords not shown take the defaults listed below).}]
Hamiltonian = DFTB {
  SCC             = Yes
  SCCTolerance    = 1e-8
  MaxSCCIterations= 400
  ...
  DeltaSCF = {
    ShiftMargin        = 0.1    # Ha; eta = |HOMO-LUMO gap| + ShiftMargin
    # UseIMOM          = Yes    # (default) greedy overlap orbital tracking
    # UseMixer         = No     # (default) direct substitution q_in := q_out
    ExcitedFrom_alpha  = {5 6}  # vacate alpha orbitals 5 and 6
    ExcitedTo_alpha    = {7 8}  # populate alpha orbitals 7 and 8
  }
  SOSCF = {
    Threshold          = 1e-4   # (default) pre-SOSCF handoff + final max|g|
    OuterSCCTolerance  = 1e-7   # charge convergence in outer loop
    # UseMixer         = No     # (default) direct substitution
    # UseMaxKappa      = Yes    # (default) trust-region step clip
    MaxKappa           = 0.2    # override (default 0.5 rad); tighter for double excitations
    Verbosity          = Yes    # detailed diagnostics in soscf.out
    ExcitedFrom_alpha  = {5 6}  # must match DeltaSCF
    ExcitedTo_alpha    = {7 8}  # must match DeltaSCF
  }
}
\end{lstlisting}

\paragraph{Keyword reference.}

\begin{table}[h]
\centering
\begin{tabular}{lll}
\toprule
Block/Keyword & Default & Meaning \\
\midrule
\multicolumn{3}{l}{\textbf{DeltaSCF block}} \\
\texttt{ShiftEnergy} & --- & Fixed level shift $\eta$ (Ha). Mutually exclusive with \texttt{ShiftMargin}. \\
\texttt{ShiftMargin} & 0.1 & $\varepsilon'$ in Eq.~\eqref{eq:stepEta}: $\eta = (\varepsilon_\mathrm{LUMO}^\mathrm{GS}-\varepsilon_\mathrm{HOMO}^\mathrm{GS})+\varepsilon'$. \\
\texttt{UseIMOM} & \texttt{Yes} & If \texttt{Yes}: IMOM in pre-SOSCF and IMOM (frozen reference) in SOSCF outer loop; if \texttt{No}: plain Aufbau in pre-SOSCF and classical MOM (reference advanced each iter) in SOSCF. \\
\texttt{UseMixer} & \texttt{No} & Broyden mixer in pre-SOSCF Delta-SCF (default: direct substitution, STEP reference). Also consulted by \texttt{runDeltaScfStandaloneLoop} when SOSCF is disabled. \\
\texttt{ExcitedFrom\_alpha} & --- & List of $\alpha$ orbital indices to vacate. \\
\texttt{ExcitedTo\_alpha}   & --- & List of $\alpha$ orbital indices to populate. \\
\texttt{ExcitedFrom\_beta}  & --- & Optional $\beta$ indices. \\
\texttt{ExcitedTo\_beta}    & --- & Optional $\beta$ indices. \\
\midrule
\multicolumn{3}{l}{\textbf{SOSCF block}} \\
\texttt{Threshold}          & $10^{-4}$ & $\vartheta_\mathrm{SOSCF}$: pre-SOSCF handoff + SOSCF gradient criterion (Ha). \\
\texttt{OuterSCCTolerance}  & $10^{-8}$  & Charge convergence criterion in the SOSCF outer loop; see Eq.~\eqref{eq:conv-SOSCF}. \\
\texttt{UseMixer}           & \texttt{No}  & Broyden/DIIS mixer inside SOSCF outer loop (default: direct substitution). \\
\texttt{UseMaxKappa}        & \texttt{Yes} & Enable trust-region clip on $\kkappa$. \\
\texttt{MaxKappa}           & 0.5 & $\kappa_\mathrm{max}$ in radians. \\
\texttt{Verbosity}          & \texttt{No} & Write detailed per-iter diagnostics to soscf.out. \\
\texttt{ExcitedFrom\_alpha}, \texttt{ExcitedTo\_alpha} & --- & Must match DeltaSCF block. \\
\texttt{ExcitedFrom\_beta}, \texttt{ExcitedTo\_beta} & --- & Optional. \\
\bottomrule
\end{tabular}
\caption{Input keywords specific to the Delta-SCF / SOSCF extension.}
\end{table}

\subsection{Data structures}

The SOSCF state is held in the derived type \texttt{TSoscf}
(\texttt{dftb/soscf.F90}):

\begin{lstlisting}[style=fortran,caption={\texttt{TSoscfSpin} (per spin channel) and \texttt{TSoscf} (top-level), from \texttt{dftb/soscf.F90}.}]
type :: TSoscfSpin
  integer :: nOcc = 0
  integer :: nVirt = 0
  integer :: nOccVirt = 0              ! = nOcc * nVirt
  integer,  allocatable :: occIdx(:), virtIdx(:)   ! 1-based indices
  real(dp), allocatable :: xVec(:)     ! accumulated rotation vector (nOccVirt)
  real(dp), allocatable :: gPrev(:)    ! previous-iter gradient (for gamma = g - gPrev)
  real(dp), allocatable :: DxPrev(:)   ! previous-iter step     (for delta = DxPrev)
  real(dp), allocatable :: invHess(:,:)! full symmetric (nOccVirt x nOccVirt) inv-Hessian
  logical :: isInitialized = .false.
end type TSoscfSpin

type :: TSoscf
  integer :: nSpin = 0
  type(TSoscfSpin), allocatable :: spin(:)
  real(dp) :: threshold   = 1.0e-4_dp   ! handoff / final max|g| threshold
  logical  :: useMaxKappa = .true.      ! enable trust-region clip
  real(dp) :: maxKappa    = 0.5_dp      ! cap in rad
end type TSoscf
\end{lstlisting}

On the run-state struct (\texttt{TDftbPlusMain} in
\texttt{initprogram.F90}) we add companion control fields and flags.
The defaults shown are the values used when the corresponding HSD
keyword is omitted:

\begin{lstlisting}[style=fortran,caption={SOSCF / Delta-SCF fields on the DFTB+ run-state struct (\texttt{initprogram.F90}).}]
! --- DeltaSCF block ---
logical  :: tDeltaScf            = .false.   ! enables pre-SOSCF or standalone Delta-SCF
logical  :: deltaScfIsFixedEta   = .false.   ! .true. if ShiftEnergy given; .false. if ShiftMargin
real(dp) :: deltaScfShiftEnergy  = 0.0_dp    ! Ha; used only when deltaScfIsFixedEta = .true.
real(dp) :: deltaScfShiftMargin  = 0.1_dp    ! Ha; epsilon' in eta = |HOMO-LUMO gap| + margin
logical  :: deltaScfUseMixer     = .false.   ! Broyden/DIIS in pre-SOSCF and standalone Delta-SCF
logical  :: deltaScfUseIMOM      = .true.    ! IMOM in pre-SOSCF; also gates IMOM vs MOM in SOSCF
integer, allocatable :: deltaScfExcitedFrom_alpha(:), deltaScfExcitedTo_alpha(:)
integer, allocatable :: deltaScfExcitedFrom_beta(:),  deltaScfExcitedTo_beta(:)

! --- SOSCF block ---
logical  :: tSoscf               = .false.   ! enables runSoscfLoop after pre-SOSCF
real(dp) :: soscfThreshold       = 1.0e-4_dp ! max|g| threshold (handoff + final)
real(dp) :: soscfOuterSccTol     = 1.0e-8_dp ! charge convergence in SOSCF outer loop
logical  :: soscfVerbose         = .false.   ! write soscf.out per-iter diagnostics
logical  :: soscfUseMixer        = .false.   ! Broyden/DIIS in SOSCF outer loop
logical  :: soscfUseMaxKappa     = .true.    ! enable trust-region clip
real(dp) :: soscfMaxKappa        = 0.5_dp    ! cap in rad
type(TSoscf) :: soscf
integer, allocatable :: soscfExcitedFrom_alpha(:), soscfExcitedTo_alpha(:)
integer, allocatable :: soscfExcitedFrom_beta(:),  soscfExcitedTo_beta(:)
\end{lstlisting}

\subsection{Key subroutines}

\paragraph{In \texttt{dftb/soscf.F90}.}
A new module encapsulating the quasi-Newton machinery:

\begin{center}
\begin{tabular}{ll}
\toprule
Routine & Purpose \\
\midrule
\texttt{soscf\_init} & Allocate per-spin arrays; build occ/virt partition from target filling; \\
& initialise diagonal inverse Hessian per Eq.~\eqref{eq:invH-init}; store $\kappa_\mathrm{max}$. \\
\texttt{computeOrbGradient} & $g_{ia} = 4\,[\CC^\top \FF\CC]_{ia}$ for given eigvecs and AO Fock. \\
\texttt{getSoscfStep} & Compute $\mathrm{D}\xx=-\HH^{-1}\gg$; apply trust-region clip if enabled. \\
\texttt{buildKappa} & Assemble antisymmetric $\kkappa$ from the flattened $\mathrm{D}\xx$. \\
\texttt{computeCayleyExp} & Cayley form $\UU=(\II-\tfrac12\kkappa)^{-1}(\II+\tfrac12\kkappa)$ via LAPACK gesv. \\
\texttt{updateInvHessianLSR1} & SR1 rank-1 update~\eqref{eq:lsr1} with two safety guards. \\
\texttt{getInvHessDiag} & Return diagonal of $\HH^{-1}$ (diagnostic). \\
\bottomrule
\end{tabular}
\end{center}

\paragraph{In \texttt{dftbplus/main.F90}.}
The driver adds:

\begin{center}
\begin{tabular}{ll}
\toprule
Routine & Purpose \\
\midrule
\texttt{runPreSoscfDeltaScfLoop} & Pre-SOSCF Delta-SCF loop; exits on $\max|g|<\vartheta_\mathrm{SOSCF}$ (Alg.~\ref{alg:deltaScf}). Called when \texttt{tSoscf = .true.}. \\
\texttt{runDeltaScfStandaloneLoop} & Standalone Delta-SCF loop (no SOSCF handoff); exits on $\mathrm{sccErrorQ}<\mathtt{SCCTolerance}$ (same criterion as GS SCC). Called when \texttt{tSoscf = .false.}. \\
\texttt{runSoscfLoop}    & SOSCF outer loop, one Newton rotation per iter (Alg.~\ref{alg:soscf}). \\
\texttt{buildDeltaScfTargetFilling} & Construct $f^{(\mathrm{T})}$ from GS filling and Excited indices. \\
\texttt{computeDeltaScfEta} & Automatic $\eta = (\varepsilon_\mathrm{LUMO}^\mathrm{GS}-\varepsilon_\mathrm{HOMO}^\mathrm{GS})+\varepsilon'$ per Eq.~\eqref{eq:stepEta}. \\
\texttt{buildAndDiagDeltaScfRealHam} & Form $\FF+\eta\,\SS\mathbf{Q}\SS$ and solve generalised eigenproblem per Eq.~\eqref{eq:step}. \\
\texttt{restoreDeltaScfEigenvalues} & Subtract $\eta$ from the eigenvalues whose filling indicates virtual ($f<1/2$). \\
\texttt{getIMOMTargetFilling} & Greedy max-overlap filling reassignment (\S\ref{sec:imom}); used both by pre-SOSCF and by \texttt{runSoscfLoop} as IMOM (frozen reference) or MOM (advancing reference). \\
\bottomrule
\end{tabular}
\end{center}

\paragraph{In \texttt{dftbplus/mainio.F90}.}
Output helpers:

\begin{center}
\begin{tabular}{ll}
\toprule
Routine & Purpose \\
\midrule
\texttt{printSoscfHeader}   & stdout 5-column header: iSCC $|$ $E_\mathrm{elec}$ $|$ $\Delta E$ $|$ SCC error $|$ $\max|g|$. \\
\texttt{printSoscfInfo}     & stdout per-iter row with the five quantities. \\
\texttt{writeSoscfHeader}   & Initialise verbose soscf.out. \\
\texttt{writeSoscfIter}     & Per-iter block in soscf.out (summary + per-orbital eigvec coefficients). \\
\texttt{writeSoscfRotDiag}  & Rotation diagnostics: L-SR1 flags, $\|\bm{\delta}\|$, $\|\bm{\gamma}\|$, denominator, $\max|\kkappa|$. \\
\texttt{writeSoscfPostRotDiag} & Post-rotation eigvectors with IMOM-assigned occupation. \\
\texttt{writeSoscfMomDiag}  & Pre/post IMOM filling and eigvec diagnostic block. \\
\bottomrule
\end{tabular}
\end{center}

\subsection{Output files}

\paragraph{stdout (the \texttt{log} file).}
Three iteration tables are printed sequentially:
\begin{enumerate}
\item Ground-state SCC (4-column: iSCC, $E_\mathrm{elec}$, $\Delta E$, SCC error).
\item Pre-SOSCF Delta-SCF (5-column with $\max|g|$ when SOSCF is active):
\begin{lstlisting}[style=hsd,numbers=none]
  Delta-SCF (pre-SOSCF) charge update: direct substitution (no mixer).
  Delta-SCF (pre-SOSCF) orbital tracking: IMOM (reference = GS eigvecs).
 iSCC Total electronic   Diff electronic      SCC error          max|g|
    1   -0.413608E+01    0.000000E+00    0.594E+00    0.000000E+00
    2   -0.353248E+01    0.603596E+00    0.092E-01    0.098371E+00
    ...
  SOSCF orbital tracking: IMOM (reference frozen at pre-SOSCF eigvecs).
\end{lstlisting}
\item SOSCF outer (same 5-column header; banner lines before are the
three lines below, each switching on the corresponding keyword —
\texttt{UseIMOM}, \texttt{UseMixer}, \texttt{UseMaxKappa}):
\begin{lstlisting}[style=hsd,numbers=none]
  SOSCF orbital tracking: IMOM (reference frozen at pre-SOSCF eigvecs).
                          # or "MOM (reference updated each outer iter)." when UseIMOM = No
  SOSCF charge update:    direct substitution (no mixer).
                          # or "Broyden/DIIS mixer." when UseMixer = Yes
  SOSCF step cap:         max|kappa| clipped to  0.500 rad (trust-region).
                          # or "disabled (full Newton step)." when UseMaxKappa = No
\end{lstlisting}
The same pair of alternatives appears in the pre-SOSCF banner
(\texttt{Delta-SCF (pre-SOSCF) orbital tracking: aufbau on
level-shifted eigenvalues.} when \texttt{UseIMOM = No}).
\end{enumerate}
The final \texttt{Total Energy} line follows the SOSCF convergence.

\paragraph{soscf.out (when \texttt{Verbosity = Yes}).}
For every pre-SOSCF and SOSCF iteration:
\begin{itemize}[leftmargin=*]
\item One summary line: iIter, $E_\mathrm{elec}$, $\Delta E$, sccErrorQ, $\max|g|$ per spin;
\item For each spin: the full eigenvalue / occupation / eigenvector-coefficient table;
\item In the SOSCF phase, additionally per spin per iter: L-SR1 flags
($\|\bm{\delta}\|, \|\bm{\gamma}\|$, denom, applied flag), max$|\kkappa|$,
full occ--virt blocks of $g$, $\mathrm{D}x$, and $\kkappa$;
\item IMOM pre/post-rotation diagnostics: \emph{Filling BEFORE IMOM},
\emph{Filling AFTER IMOM}, reference and candidate eigvectors.
\end{itemize}

\paragraph{band.out.}
Eigenvalues and fillings of the final converged state. After a SOSCF
run \texttt{band.out} is written \emph{after} the post-SOSCF eigenvalue
refresh $\varepsilon_i \gets \langle\CC_i|\FF|\CC_i\rangle$, so the
reported values correspond self-consistently to the rotated MOs in
\texttt{eigvecsReal}. The $(\varepsilon_i,f_i)$ pairs in
\texttt{band.out} are additionally sorted ascending by $\varepsilon$
(the internal arrays remain in IMOM order). For a standalone
Delta-SCF run \texttt{band.out} is written from the converged
level-shifted eigenvalues directly, without the SOSCF refresh and
without sorting.

\paragraph{detailed.out, charges.bin.}
Standard DFTB+ outputs, using (in SOSCF mode) the refreshed eigenvalues
and IMOM-tracked filling in the internal (unsorted) order.

\subsection{Convergence criteria}

Let $\vartheta_\mathrm{g}\equiv\vartheta_\mathrm{SOSCF}=\mathtt{SOSCF\,\{\,Threshold\,\}}$
and $\mathrm{sccErrorQ}\equiv\max_i|q^\mathrm{out}_i-q^\mathrm{in}_i|$.
Three exit rules are in play, one per loop:

\begin{equation}
\text{Pre-SOSCF (\texttt{runPreSoscfDeltaScfLoop}) exit:}\quad
\max_{ia}|g_{ia}| \,<\, \vartheta_\mathrm{g}.
\label{eq:conv-preSOSCF}
\end{equation}
$\mathrm{sccErrorQ}$ is evaluated and printed per iter for diagnostics
but is \emph{not} used as an exit condition here — charge
self-consistency is the SOSCF outer loop's job.

\begin{equation}
\text{SOSCF (\texttt{runSoscfLoop}) converged:}\quad
\mathrm{sccErrorQ}\,<\,\mathtt{OuterSCCTolerance}
\;\text{\textbf{and}}\;
\max_{ia}|g_{ia}| \,<\, \vartheta_\mathrm{g}.
\label{eq:conv-SOSCF}
\end{equation}
Both the orbital gradient (stationary point on the rotation manifold)
and the charge residual (self-consistent density at that stationary
point) must be small simultaneously. The two quantities are now
evaluated at \emph{different} states of the current outer iteration
$n$, by design: $\max|\gg|$ is measured \emph{before} the rotation at
$(\CC^{(n-1)},\FF[\qq^{(n-1),\mathrm{in}}])$ in Step~A, whereas
$\mathrm{sccErrorQ}=\|\qq^\mathrm{out}(\CC^{(n)})-\qq^{(n-1),\mathrm{in}}\|_\infty$
is measured \emph{after} the rotation in Step~B (post-rotation $\CC$
against the $\qq^\mathrm{in}$ that built the Fock matrix used in
Step~A). At convergence the last rotation has
$\max|\kkappa|=O(\max|\gg|/\|\HH\|)$, so the one-rotation lag on the
gradient is $O(\vartheta_\mathrm{SOSCF})$ --- below the convergence
noise; see \S\ref{sec:algo-soscf}.

\begin{equation}
\text{Standalone (\texttt{runDeltaScfStandaloneLoop}) converged:}\quad
\mathrm{sccErrorQ}\,<\,\mathtt{SCCTolerance}
\;\;\text{(global SCC tolerance).}
\label{eq:conv-standalone}
\end{equation}
This reproduces the behaviour of the ground-state SCC loop applied to
the level-shifted fixed point; no gradient is computed.

\subsection{Results on selected test cases}\label{sec:results-dblexc}

Table~\ref{tab:results} summarises behaviour on three excitations of
\ce{H2O} and formaldehyde:

\begin{table}[ht]
\centering
\small
\setlength{\tabcolsep}{5pt}
\begin{tabular}{@{}lrrr@{}}
\toprule
Case & pre-SOSCF iters & SOSCF iters & $E_\mathrm{tot}$ (Ha) \\
\midrule
\ce{H2O}, HOMO$\to$LUMO                                     & 121 & 118 & $-3.4802356738$ \\
\ce{H2CO}, HOMO-1$\to$LUMO                                  & 5   & 94  & $-5.4813054440$ \\
\ce{H2CO}, $\{$HOMO-1,\,HOMO$\}\!\to\!\{$LUMO,\,LUMO+1$\}$  & 2   & 328 & $-4.7979030943$ \\
\bottomrule
\end{tabular}
\caption{All cases use the current defaults (\texttt{UseIMOM=Yes},
\texttt{UseMixer=No} for both blocks, \texttt{UseMaxKappa=Yes},
\texttt{MaxKappa=0.5}~rad). The double excitation requires
\texttt{MaxKappa=0.2} to converge; without the trust-region clip it
diverges exponentially ($\max|g|>10^{77}$).}
\label{tab:results}
\end{table}

The $\kappa_\mathrm{max}$ default of 0.5~rad is inactive for the first
two cases (natural Newton step is well below 0.5~rad once SOSCF is in
its convergence regime). Double excitations and other high-order
saddles may require tightening to 0.2~rad.

%=====================================================================
\section{Summary}

The implementation delivers a second-order SCF method for excited
states in DFTB+ built around four interacting ingredients:
\begin{enumerate}[leftmargin=*]
\item Level-shift Delta-SCF (STEP) for a warm start to the excited
saddle;
\item IMOM greedy-overlap orbital tracking that stabilises non-aufbau
character under SCC drift;
\item Quasi-Newton step via L-SR1 inverse-Hessian update with the
correct saddle-seeking signed Newton direction;
\item Cayley-form orbital rotation \eqref{eq:cayley} guaranteeing
orthonormality, together with a trust-region cap
\eqref{eq:trust-cap} that prevents catastrophic overshoot on high-order
saddles.
\end{enumerate}
The input is controlled by the \texttt{DeltaSCF\{\dots\}} and
\texttt{SOSCF\{\dots\}} blocks inside the \texttt{Hamiltonian=DFTB} HSD
block; outputs include an extended stdout iteration log with a
$\max|g|$ column, the detailed per-iter diagnostic file \texttt{soscf.out},
and a refreshed \texttt{band.out} self-consistent with the post-rotation
MOs.

\subsection*{Full pseudocode}
\addcontentsline{toc}{subsection}{Full pseudocode}

Algorithm~\ref{alg:full} below lists the complete driver at one level
of detail, making explicit what each step updates. Throughout the
listing, $\qq$ denotes the (reduced) charge vector,
$\CC$ the MO coefficients, $\bm{\varepsilon}$ the eigenvalues,
$f$ the occupation array, $\FF$ the AO Fock/Hamiltonian,
$\SS$ the AO overlap, $\gg$ the orbital gradient,
$\HH^{-1}$ the approximate inverse orbital Hessian, and
$\mathrm{D}\xx$ the quasi-Newton step in the flattened occ--virt
parameter space.

The full driver is shown in Algorithms~\ref{alg:full-init}--\ref{alg:full-SOSCF},
one algorithm per phase group to keep each block on a single page.

\medskip

\textbf{Input.} Geometry, basis, user-specified
\texttt{ExcitedFrom} / \texttt{ExcitedTo} indices, level-shift parameter
($\eta$ or $\varepsilon'$), thresholds $\mathtt{SCCTolerance}$,
$\vartheta_\mathrm{SOSCF}$, $\mathtt{OuterSCCTolerance}$, flags
\texttt{tSoscf}, \texttt{UseIMOM}, \texttt{UseMixer}, \texttt{UseMaxKappa}.\\
\textbf{Output.} Converged $(\CC, \bm{\varepsilon}, f, \qq, E)$
and files \texttt{band.out} / \texttt{detailed.out} / \texttt{charges.bin}.

\begin{algorithm}[H]
\caption{Phases 0--1: ground-state SCC and target-filling setup.}
\label{alg:full-init}
\footnotesize
\begin{algorithmic}[1]

\State \textbf{Phase 0 --- Ground-state SCC.}
\State Initialise $\qq^{(0)}\gets\qq^{0}$ (neutral-atom reference).
\For{$k=1,2,\dots,\mathrm{MaxSCC}$}
  \State Build $\FF[\qq^{(k-1)}]$ from Eq.~\eqref{eq:dftbH}.
  \State Solve $\FF\,\CC=\SS\,\CC\,\bm{\varepsilon}$ $\to$ $(\CC,\bm{\varepsilon})$.
  \State Apply Aufbau $\to$ $f$; form $\mathbf{P}$; compute $\qq^{\mathrm{out}}$ (Eq.~\eqref{eq:mulliken}).
  \If{$\|\qq^{\mathrm{out}}-\qq^{(k-1)}\|_\infty<\mathtt{SCCTolerance}$}\ \textbf{break} \EndIf
  \State $\qq^{(k)}\gets\mathrm{Mixer}(\qq^{(k-1)},\qq^{\mathrm{out}}-\qq^{(k-1)})$ \Comment{Broyden/DIIS}
\EndFor
\State Save GS state: $(\CC^{(\mathrm{GS})},\bm{\varepsilon}^{(\mathrm{GS})},\qq^{(\mathrm{GS})})\gets(\CC,\bm{\varepsilon},\qq)$.

\Statex

\State \textbf{Phase 1 --- Build non-Aufbau target filling.}
\State $f^{(\mathrm{T})} \gets f^{(\mathrm{GS})}$; set $f^{(\mathrm{T})}_r=0$ for $r\in\mathtt{ExcitedFrom}$, $f^{(\mathrm{T})}_r=1$ for $r\in\mathtt{ExcitedTo}$ (per spin).
\State Determine level shift once: $\eta\gets\mathtt{ShiftEnergy}$,
       \textbf{or} $\eta = \varepsilon_\mathrm{LUMO}^{(\mathrm{GS})}-\varepsilon_\mathrm{HOMO}^{(\mathrm{GS})}+\varepsilon'$.

\end{algorithmic}
\end{algorithm}

\begin{algorithm}[H]
\caption{Phase 2: level-shifted Delta-SCF. Exactly one of 2a or 2b runs,
selected by \texttt{tSoscf}.}
\label{alg:full-deltaSCF}
\footnotesize
\begin{algorithmic}[1]

\If{\textbf{not} \texttt{tSoscf}} \Comment{Standalone Delta-SCF branch}
  \State \textbf{Phase 2a --- runDeltaScfStandaloneLoop.}
  \If{\texttt{UseIMOM}} $\CC^{(\mathrm{ref})}\gets\CC^{(\mathrm{GS})}$;\ $f^{(\mathrm{ref})}\gets f^{(\mathrm{T})}$. \EndIf
  \State $f\gets f^{(\mathrm{T})}$;\ $\qq^{(0)}\gets\qq^{(\mathrm{GS})}$.
  \For{$k=1,2,\dots,\mathrm{MaxSCC}$}
    \State \textbf{[shift]} Form $\mathbf{Q}=\sum_{a\in\mathrm{virt}(f)}\CC_{:,a}\CC_{:,a}^{\top}$; solve $(\FF[\qq^{(k-1)}]+\eta\,\SS\mathbf{Q}\SS)\CC=\SS\,\CC\,\widehat{\bm{\varepsilon}}$.
    \State \textbf{[fill]} Aufbau on $\widehat{\bm{\varepsilon}}$ $\to$ $f$; \textbf{if} \texttt{UseIMOM} reassign $f$ by max-overlap.
    \State \textbf{[restore]} $\varepsilon_a\gets\widehat{\varepsilon}_a-\eta$ for $a\in\mathrm{virt}(f)$.
    \State \textbf{[charges]} Build $\mathbf{P}$; $\qq^{\mathrm{out}}\gets\mathrm{Mulliken}$; $\mathrm{sccErrorQ}\gets\|\qq^{\mathrm{out}}-\qq^{(k-1)}\|_\infty$.
    \If{$\mathrm{sccErrorQ}<\mathtt{SCCTolerance}$}\ \textbf{break (done)} \EndIf
    \State $\qq^{(k)}\gets\mathrm{Mixer}(\qq^{(k-1)},\qq^{\mathrm{out}}-\qq^{(k-1)})$ if \texttt{UseMixer} \textbf{else} $\qq^{(k)}\gets\qq^{\mathrm{out}}$.
  \EndFor
  \State \textbf{goto Phase 4} (skip SOSCF).
\Else

  \Statex
  \State \textbf{Phase 2b --- runPreSoscfDeltaScfLoop.} \Comment{warm start for SOSCF}
  \If{\texttt{UseIMOM}} $\CC^{(\mathrm{ref})}\gets\CC^{(\mathrm{GS})}$;\ $f^{(\mathrm{ref})}\gets f^{(\mathrm{T})}$. \EndIf
  \State $f\gets f^{(\mathrm{T})}$;\ $\qq^{(0)}\gets\qq^{(\mathrm{GS})}$.
  \For{$k=1,2,\dots,\mathrm{MaxSCC}$}
    \State \textbf{[shift]} solve $(\FF[\qq^{(k-1)}]+\eta\,\SS\mathbf{Q}\SS)\CC=\SS\,\CC\,\widehat{\bm{\varepsilon}}$ \ ($\mathbf{Q}$ from current $f$).
    \State \textbf{[fill]} Aufbau on $\widehat{\bm{\varepsilon}}$ $\to$ $f$; \textbf{if} \texttt{UseIMOM} reassign $f$ by max-overlap.
    \State \textbf{[restore]} $\varepsilon_a\gets\widehat{\varepsilon}_a-\eta$ for $a\in\mathrm{virt}(f)$.
    \State \textbf{[gradient]} $g_{ia}\gets 4\,[\CC^{\top}\FF[\qq^{(k-1)}]\CC]_{ia}$ on occ--virt.
    \If{$\max_{ia}|g_{ia}|<\vartheta_\mathrm{SOSCF}$}\ \textbf{break (handoff)} \EndIf
    \State \textbf{[charges]} Build $\mathbf{P}$; $\qq^{\mathrm{out}}\gets\mathrm{Mulliken}$; compute $\mathrm{sccErrorQ}$ (diagnostic only).
    \State $\qq^{(k)}\gets\mathrm{Mixer}(\qq^{(k-1)},\qq^{\mathrm{out}}-\qq^{(k-1)})$ if \texttt{UseMixer} \textbf{else} $\qq^{(k)}\gets\qq^{\mathrm{out}}$.
  \EndFor
  \State $f^{(\mathrm{T})}\gets f$ \Comment{final IMOM-tracked filling handed to SOSCF}

\EndIf

\end{algorithmic}
\end{algorithm}

\begin{algorithm}[H]
\caption{Phase 3--4: SOSCF outer loop (skipped if Phase 2a ran) and
output.}
\label{alg:full-SOSCF}
\footnotesize
\begin{algorithmic}[1]

\State \textbf{Phase 3 --- runSoscfLoop.}
\State \texttt{soscf\_init}: partition occ/virt via $f^{(\mathrm{T})}$; set
$\HH^{-1}\gets\mathrm{diag}\!\big(\mathrm{sign}(\Delta_{ia})/\max(2|\Delta_{ia}|,10^{-3})\big)$
with $\Delta_{ia}=\varepsilon_a-\varepsilon_i$ (Eq.~\eqref{eq:invH-init}).
\State Initialise tracking reference: $\CC^{(\mathrm{ref})}\gets\CC$,\ $f^{(\mathrm{ref})}\gets f^{(\mathrm{T})}$.
\State Zero history: $\gg_{\mathrm{prev}}\gets 0$,\ $\mathrm{D}\xx_{\mathrm{prev}}\gets 0$.
\State Reset Broyden/DIIS mixer once (if \texttt{UseMixer}). $\qq^{(0)}\gets$ charges from end of Phase 2b.
\For{$n=1,2,\dots,\mathrm{MaxSCC}$}
  \Statex\quad\textit{Step A --- one Newton rotation at rebuilt $\FF[\qq^{(n-1)}]$.}
  \State Rebuild $\FF\gets\FF[\qq^{(n-1)}]$ from the current input charges.
  \For{each spin channel}
    \State $\gg\gets 4\,[\CC^{\top}\FF\,\CC]_{ia}$.
    \If{$n>1$}
      \State $\bm{\gamma}\gets\gg-\gg_{\mathrm{prev}}$,\ $\bm{\delta}\gets\mathrm{D}\xx_{\mathrm{prev}}$,\ $\bm{j}\gets\bm{\delta}-\HH^{-1}\bm{\gamma}$.
      \If{safeguards pass} $\HH^{-1}\gets\HH^{-1}+\bm{j}\bm{j}^{\top}/(\bm{j}^{\top}\bm{\gamma})$ \EndIf\Comment{Eq.~\eqref{eq:lsr1}}
    \EndIf
    \State $\gg_{\mathrm{prev}}\gets\gg$; \ $\mathrm{D}\xx\gets-\HH^{-1}\gg$; clip to $\kappa_\mathrm{max}$ if \texttt{UseMaxKappa}; \ $\mathrm{D}\xx_{\mathrm{prev}}\gets\mathrm{D}\xx$.
    \State Build antisymmetric $\kkappa$; form $\UU=(\II-\tfrac12\kkappa)^{-1}(\II+\tfrac12\kkappa)$; rotate $\CC\gets\CC\,\UU$.
    \State Reassign $f^{(\mathrm{T})}$ by max-overlap of $\CC$ against $(\CC^{(\mathrm{ref})},f^{(\mathrm{ref})})$.
  \EndFor
  \If{\textbf{not} \texttt{UseIMOM}} $\CC^{(\mathrm{ref})}\gets\CC$,\ $f^{(\mathrm{ref})}\gets f^{(\mathrm{T})}$. \EndIf \Comment{MOM reference advance}

  \Statex\quad\textit{Step B --- density, charges and energies at the post-rotation $\CC$.}
  \State Build $\mathbf{P}$ from $\CC$ and $f^{(\mathrm{T})}$; $\qq^{\mathrm{out}}\gets\mathrm{Mulliken}$.
  \State Evaluate total electronic energy $E_\mathrm{elec}$ at the rotated state (\texttt{calcEnergies}, \texttt{sumEnergies}).
  \State $\mathrm{sccErrorQ}\gets\|\qq^{\mathrm{out}}-\qq^{(n-1)}\|_\infty$.
  \State $\qq^{(n)}\gets\mathrm{Mixer}(\qq^{(n-1)},\qq^{\mathrm{out}}-\qq^{(n-1)})$ if \texttt{UseMixer} \textbf{else} $\qq^{(n)}\gets\qq^{\mathrm{out}}$. \Comment{feeds Step~A of iter $n{+}1$}

  \Statex\quad\textit{Step C --- convergence check.}
  \If{$\mathrm{sccErrorQ}<\mathtt{OuterSCCTolerance}$ \textbf{and} $\max|\gg|<\vartheta_\mathrm{SOSCF}$}\ \textbf{break} \EndIf
\EndFor
\State \textbf{Refresh eigenvalues:} $\varepsilon_i\gets\langle\CC_i|\FF|\CC_i\rangle$.

\Statex

\State \textbf{Phase 4 --- Output.} Write \texttt{band.out} (sorted by $\varepsilon$ after SOSCF; unsorted otherwise), \texttt{detailed.out}, \texttt{charges.bin} from $(\CC,\bm{\varepsilon},f^{(\mathrm{T})},\qq)$.
\end{algorithmic}
\end{algorithm}

The three principal state variables evolving across the driver are
therefore:
\begin{itemize}[leftmargin=*,nosep]
\item \textbf{$\qq$} — updated once per SCC/SOSCF outer iteration by
Mulliken analysis plus (optional) Broyden/DIIS mixing; carries the
self-consistent density.
\item \textbf{$(\CC,f)$} — updated by \emph{diagonalisation} of the
level-shifted $\FF$ during GS SCC and Delta-SCF phases, and by
\emph{orbital rotation} $\CC\gets\CC\,\UU$ inside the SOSCF outer
loop; the filling $f$ is set by Aufbau and then (optionally) overridden
by IMOM/MOM overlap matching.
\item \textbf{$\HH^{-1}$} — initialised diagonally from the signed
occ--virt eigenvalue gaps and refined by the L-SR1 rank-1 update once
per SOSCF outer iteration.
\end{itemize}

%=====================================================================
\begin{thebibliography}{9}

\bibitem{FischerAlmlof1992}
T.\ H.\ Fischer and J.\ Alml\"of,
``General Methods for Geometry and Wave Function Optimization,''
\emph{J.\ Phys.\ Chem.} \textbf{96}, 9768 (1992).

\bibitem{Neese2000}
F.\ Neese,
``Approximate second-order SCF convergence for spin unrestricted
wavefunctions,''
\emph{Chem.\ Phys.\ Lett.} \textbf{325}, 93 (2000).

\bibitem{Bacskay1981}
G.\ B.\ Bacskay,
``A quadratically convergent Hartree-Fock (QC-SCF) method.
Application to closed shell systems,''
\emph{Chem.\ Phys.} \textbf{61}, 385 (1981).

\bibitem{HelgakerBook}
T.\ Helgaker, P.\ J\o{}rgensen, J.\ Olsen,
\emph{Molecular Electronic Structure Theory},
Wiley, Chichester (2000).

\bibitem{CarterFenkHerbert2020}
K.\ Carter-Fenk and J.\ M.\ Herbert,
``State-Targeted Energy Projection: A Simple and Robust Approach to
Orbital Relaxation of Non-Aufbau Self-Consistent Field Solutions,''
\emph{J.\ Chem.\ Theory Comput.} \textbf{16}, 5067 (2020).

\bibitem{GilbertBesleyGill2008}
A.\ T.\ B.\ Gilbert, N.\ A.\ Besley, P.\ M.\ W.\ Gill,
``Self-consistent field calculations of excited states using the
maximum overlap method (MOM),''
\emph{J.\ Phys.\ Chem.\ A} \textbf{112}, 13164 (2008).

\bibitem{Barca2018}
G.\ M.\ J.\ Barca, A.\ T.\ B.\ Gilbert, P.\ M.\ W.\ Gill,
``Simple Models for Difficult Electronic Excitations,''
\emph{J.\ Chem.\ Theory Comput.} \textbf{14}, 1501 (2018).

\bibitem{HaitHeadGordon2020}
D.\ Hait and M.\ Head-Gordon,
``Excited State Orbital Optimization via Minimizing the Square of the
Gradient: General Approach and Application to Singly and Doubly
Excited States,''
\emph{J.\ Chem.\ Theory Comput.} \textbf{16}, 1699 (2020).

\bibitem{NocedalWright}
J.\ Nocedal and S.\ J.\ Wright,
\emph{Numerical Optimization}, 2nd ed.,
Springer, New York (2006).

\bibitem{ConnGouldToint}
A.\ R.\ Conn, N.\ I.\ M.\ Gould, P.\ L.\ Toint,
\emph{Trust-Region Methods},
MPS-SIAM, Philadelphia (2000).

\bibitem{Elstner1998}
M.\ Elstner, D.\ Porezag, G.\ Jungnickel, J.\ Elsner, M.\ Haugk, Th.\ Frauenheim,
S.\ Suhai, G.\ Seifert,
``Self-consistent-charge density-functional tight-binding method for simulations of
complex materials properties,''
\emph{Phys.\ Rev.\ B} \textbf{58}, 7260 (1998).

\bibitem{DFTBplus}
B.\ Hourahine \emph{et al.},
``DFTB+, a software package for efficient approximate density
functional theory based atomistic simulations,''
\emph{J.\ Chem.\ Phys.} \textbf{152}, 124101 (2020).

\end{thebibliography}

\end{document}
